% \section{Software Infrastructure}

\section{Software Implementation}
\label{sec:softwareinfrastructure}
\label{sec:softwarecomponents}

\textbf{Security Monitor}
The SM does not allocate any resources, this task is left to the OS, but it interposes between each security-sensitive operation (like a resource allocation) and checks that the operating system performed the task without violating the high-level security policies (e.g., it is allocating a memory region to an existing enclave).
As a result, the SM is mostly responsible for bookkeeping operations, performing security checks and setting up the hardware mechanisms to enforce isolation between the running security domains.
It will not only systematically intervene on any security-sensitive operations, but will also play the role of a small kernel and link the OS to the hardware on some non-security-sensitive operations, like sending inter-process interrupts (IPI).
A more detailed description of the SM data structures and API, can be found in \cite{lebedev2019sanctorum}.

\textbf{The SM Kernel Module}
The main SM should not be possibly made unavailable to the OS by a user-level program.
Furthermore, no resource should be allocated to a user-level program without OS approval.
As a result, to create an enclave or perform other operations,
a user-application needs to interact with the SM through the SM Kernel Module.
Because the SM Kernel Module is part of the OS, it can refuse requests if needed.

\textbf{Mini-SM}
Because the SM is a shared resource, its sharing might leak some sensitive information. For instance, an enclave might call the SM at particular times and in a pattern that might leak sensitive information regarding the enclave's secrets.
As a result, each enclave uses its own private copy of the SM.
This "mini-SM" is smaller than the one accessible to the OS and only contains the code corresponding to the subset of the SM API accessible by the enclave.
It also uses its own private stack.
This mini-SM is located in private enclave memory (not accessible nor visible by the OS) and protected using physical memory isolation (hence is neither modifiable nor accessible by the enclave).
The enclave interacts with its mini-SM through the ecall API.
Note that some of these ecalls might touch shared global SM state and locks.
As a result, an enclave should be careful when interacting with the SM and be aware that it might still leak timing information to other programs.
We emphasize that an enclave binary does not go through the SM Kernel Module and directly interacts with its mini-SM through ecalls.

\textbf{Companion User-Level Application}
The companion user-level application runs on top of the OS and requests to start the enclave, and interacts with the enclave.
It does not benefit from any integrity or privacy guarantees, unlike its enclave counterpart.
However, it has access to a rich interface with the outside world (like the file system or the network) through all the usual OS services implemented through system calls.
Even if not required for an enclave to work, the companion user-level application is a good ally to an enclave program.
It can, for instance, fetch encrypted private data and forward it to the enclave or perform many other non-security-sensitive tasks.
As a result, it is important for an enclave to have a rich interface with its user program counterpart using shared memory (see Section \ref{sec:sharedmemory}).